\chapter{Part 4 Title}

% Replace this example text with the content of Part 4.
\section{Overview}
Write a brief introduction to this part and explain how it relates to the overall dossier.

\section{Main discussion}
Organize the material with sections and subsections as needed. You can include
figures, tables, equations, and citations here.

% What are x-rays (e.g., in terms of wavelength and type of radiation)
X-rays are highly-energetic wavelets of electromagnetic radiation (light). The wavelength is between ultraviolet light and gamma-rays, roughly in the range of $10^{-10\pm 2}$. The attenuation of x-rays is dependent on the material. This can be used to work out the type of material. 

% How are x-rays produced and what are the key parameters defining the nature of an x-ray source and how do different sources differ in their mechanisms and parameters
1. Lab source
	easy to use but poor resolution
	
2. Synchrotron
	better resolution but restricted availability
	

	
X-rays can be produced by exciting electrons, and allowing them to relax thereby producing a photon with a predictable energy (given by the difference in energy level between the states).

X-rays can also be produced by accelerating electrons (using for instance a linear accelerator), and then changing their trajectory by forcing them around a circle, which produces light. The electrons are accelerated into the storage storage ring, which are steered into a circle by the use of strong magnets.

They can also be accelerated into a target, which produces light via the nuclear interaction between the electron and target. The key parameters which determine the energy of the photons is the voltage of the 

% How do x-rays interact with matter and which are the important mechanisms for experimental mechanics? What makes x-rays a good probe for experimental mechanics?
The higher the energy, the larger the cross-section.

The convergance of the beam is very important for the detail of the imaging. If the beam is very convergent, we can infer from the divergance of the image some information about the specimen.

High flux gives fast imaging.

% If you wanted to study hard matter (e.g., metals) what might you need in terms of the x-ray energy? (and conversely what would you need for soft matter, e.g., polymers?)

Wavelength is comparable to distance between crystals, so we can use diffraction.

High-photon energy ionizes atoms, so they modify the material (destructive). 

X-rays interact both magnetically and electrically, but the magnetic interaction is very weak.

There are five typical interactions, of which the first three in the list are the most intersting for x-ray imaging in mechanics.

Incoherent scattering not too interesting for us, as it's random? What does that mean.

Coherent scattering is useful since its elastic, and we can predict the result.

